\section{Supplementary experimental details}
\label{app:records}

\subsection{Material availability}
An NDA restricts disclosure of the pipeline code and its implementation
components. We are seeking permission to publish further materials;
at present, we can provide execution logs.

\subsection{Execution trace example}
We illustrate the recorded decisions with BIRD development question~1068,
from the May~27 run using Claude Opus~4.7. The question asks:
``From 2010 to 2015, what was the average overall rating of players who are
higher than 170?'' It uses the football database.
Listing~\ref{lst:trace-query} shows the SQL before the division repair.

\begin{lstlisting}[language=SQL,caption={Generated SQL before the division repair.},label={lst:trace-query},captionpos=t]
SELECT SUM(t2.overall_rating) / COUNT(t2.id)
FROM Player AS t1
INNER JOIN Player_Attributes AS t2
  ON t1.player_api_id = t2.player_api_id
WHERE strftime('%Y', t2.date) >= '2010'
  AND strftime('%Y', t2.date) <= '2015'
  AND t1.height > 170
\end{lstlisting}

The execution log reports one returned row, a judge score of 90, and the
decision to keep the query. The later deterministic step detects division
without an explicit real-valued operand and inserts a cast. The following
log excerpt omits timestamps and logger prefixes:

\begin{lstlisting}[caption={Recorded execution and repair decisions.},label={lst:trace-decisions},captionpos=t]
[exec:done] 0.26s rows=1
[judge] confidence=90 (split)
[judge] confidence >= 90 - keeping original
[cast-real-div] wrapped 1 numerator(s)
  in CAST AS REAL
\end{lstlisting}

This development version used a judge threshold of 90; the final
configuration is described in Section~\ref{sec:pipeline}. The repaired
projection is:

\begin{lstlisting}[language=SQL]
CAST(SUM(t2.overall_rating) AS REAL)
  / COUNT(t2.id)
\end{lstlisting}

Replaying the saved input SQL returns 68. The repaired query returns
68.89975315514602. This matches the archived reference result. Thus the
recorded rewrite changes this answer from incorrect to correct. The value
before repair was obtained by SQL replay; the final value and verdict were
stored in the prediction record.

The question identifier links the prediction to its step records. Phoenix
spans retain step inputs and outputs, timing, and model-call usage where
available. These records connect the final verdict to the SQL change that
produced it, even when an earlier judge accepted the original query.

\subsection{Inference usage and timing}
\label{app:usage}
Table~\ref{tab:archive-cost} reports the June~4 full-pipeline run on
1,534 BIRD development questions with Claude Opus~4.8.

\begin{table}[ht]
\caption{Model usage and question-processing times, June 4. Input tokens
include cache reads and writes.}
\label{tab:archive-cost}
\centering\footnotesize
\begin{tabular}{@{}lr@{}}
\toprule
Measure & Value \\
\midrule
Input tokens & 184,173,469 \\
Output tokens & 937,571 \\
Model calls & 11,310 \\
Calls per question & 7.37 \\
Median duration (s) & 18.69 \\
95th percentile (s) & 36.80 \\
\bottomrule
\end{tabular}
\end{table}

The call count includes 22 failed calls. Durations use 1,533 completed
question-processing spans; four unfinished attempts are excluded and one
question's root span is missing. The 95th percentile uses the nearest-rank
method. Both model slots used Claude Opus~4.8 with a maximum output length
of 4,096 tokens and no explicit temperature setting. Token totals depend
on model configuration and cache behavior.

\subsection{Complete log for one question}
The listing contains every recorded line for question~1068 in the May~27
run, from its benchmark entry to the start of the next question.
Timestamps, logger names, and SQL are retained without omissions.
The complete 1,534-question run log is supplied separately as
\texttt{assets/logs/run-2026-05-27.log}.

\lstinputlisting[caption={Complete recorded log for question 1068.},
label={lst:debug-log},captionpos=t,
escapeinside={(*@}{@*)}
]{assets/logs/question-1068-print.txt}
